
\section{Rigorous quadrature for residual norms}\label{sec:applications}

In this section we return to a posteriori verification of neural network-based PDE solvers. We first provide concrete residual energy estimates on the form \eqref{eq:heat-wave-model-framework}, and discuss how these residual norms will be rigorously bounded using the quadrature framework described in \cref{subsec:quadrature}. Then we detail how the derivative bounds from \cref{sec:activation} are used to compute the local error bound \eqref{eq:cell-model-errors} for each residual norm.  Finally, we present numerical results for the heat and wave equations.



\subsection{Explicit energy estimates}\label{sec:energy-estimates}

We first derive explicit energy estimates for the heat and wave equations, which will be used to establish the residual framework \eqref{eq:heat-wave-model-framework}.

\begin{notation}\label{not:heat-wave-operator}
Throughout, \(H_0^1(\Omega)\) is equipped with the norm
\[
    \norm{v}_{H_0^1(\Omega)}\coloneqq\norm{\nabla v}_{L^2(\Omega)},
\]
and \(H^{-1}(\Omega)\) denotes the dual space with respect to this norm.
Let \(A=-\Delta\) be the positive self-adjoint operator on \(L^2(\Omega)\)
associated with the form
\[
    a(v,w)=\int_\Omega \nabla v\cdot\nabla w\,dx,
    \qquad v,w\in H_0^1(\Omega).
\]
Let \((w_k,\lambda_k)_{k\ge1}\) denote
the eigenpairs of \(A\) which form an orthonormal basis of
\(L^2(\Omega)\) and are orthogonal in \(H_0^1(\Omega)\)
\cite[Section 6.5, Theorem 1]{evans_2010}, with
\begin{equation}\label{eq:dirichlet-eigenbasis}
    Aw_k=\lambda_k w_k,
    \qquad
    0<\lambda_1\le\lambda_2\le\cdots .
\end{equation}
For \(m\ge1\), write
\(
    V_m\coloneqq \operatorname{span}\{w_1,\dots,w_m\},
\)
and, whenever Galerkin approximations are used, take \(u_m(t)\in V_m\) to
satisfy the corresponding weak form for all test functions in \(V_m\) and a.e.
\(t\in(0,T)\).

\end{notation}


We then have the explicit Poincare inequality
\begin{equation}\label{eq:poincare-h01}
    \norm{v}_{L^2(\Omega)}
    \le
    \lambda_1^{-1/2}\norm{v}_{H_0^1(\Omega)},
    \qquad v\in H_0^1(\Omega).
\end{equation}


\begin{lemma}\label{lem:convex-domain-A-H2}
Assume that \(\Omega\) is convex. Then every \(w\in D(A)\) belongs
to \(H^2(\Omega)\cap H_0^1(\Omega)\), and
\[
    \norm{w}_{H^2(\Omega)}
    \le
    C_\Omega\norm{Aw}_{L^2(\Omega)},
    \qquad
    C_\Omega
    =
    \left(
        1+\frac1{\lambda_1}+\frac1{\lambda_1^2}
    \right)^{1/2}.
\]
\end{lemma}

\begin{proof}
Since \(w\in D(A)\), we have \(Aw\in L^2(\Omega)\) and \(w\) solves the
Dirichlet problem \(-\Delta w=Aw\) weakly.  Convex-domain elliptic regularity
\cite[Theorem 3.2.1.2]{grisvard_1985} gives
\(w\in H^2(\Omega)\cap H_0^1(\Omega)\).  Then it is a classical result that
\[
    \norm{D^2 w}_{L^2(\Omega)}
    \le
    \norm{Aw}_{L^2(\Omega)}.
\]
Moreover, integration by parts and Poincare's inequality give
\[
    \norm{\nabla w}_{L^2}
    \le
    \lambda_1^{-1/2}\norm{Aw}_{L^2},
    \qquad
    \norm{w}_{L^2}
    \le
    \lambda_1^{-1}\norm{Aw}_{L^2}.
\]
Combining these three inequalities gives the claim.
\end{proof}

\subsubsection{General heat equation energy estimates}

\begin{theorem}\label{thm:heat-estimates}
Let \(\Omega\) be a open, bounded set and $\kappa, T > 0$.
Consider the heat equation
\begin{equation}\label{eq:heat-general}
\begin{cases}
\partial_t u-\kappa\Delta u=F & \text{in }\Omega\times(0,T],\\
u=0 & \text{on }\partial\Omega\times[0,T],\\
u=g & \text{on }\Omega\times\set{t=0},
\end{cases}
\end{equation}
where \(F\in L^2(0,T;L^2(\Omega))\) and \(g\in L^2(\Omega)\).  There is a
unique weak solution
\[
    u\in C([0,T];L^2(\Omega))
    \qquad
    u'\in L^2(0,T;H^{-1}(\Omega)),
\]
for which the following is true:
\begin{enumerate}[label=\textup{(\arabic*)}]
\item\label{thm:heat-estimates-l2} $u \in L^2(0,T;H_0^1(\Omega))$ with the bound
\begin{equation}\label{eq:explicit-L2-estimate}
\begin{split}
    &\sup_{0\le t\le T}\norm{u(t)}_{L^2(\Omega)}
    + \norm{u}_{L^2(0,T;H_0^1(\Omega))}
    + \norm{u'}_{L^2(0,T;H^{-1}(\Omega))}
    \\
    &\qquad\le
    \alpha^{\mathrm H}\norm{g}_{L^2(\Omega)}
    +
    \beta^{\mathrm H}
    \norm{F}_{L^2(0,T;L^2(\Omega))},
\end{split}
\end{equation}
where
\begin{equation}\label{eq:heat-alpha-beta-constants}
    \alpha^{\mathrm H}\coloneqq 1+\frac1{\sqrt{\kappa}}+\sqrt{\kappa},
    \qquad
    \beta^{\mathrm H}\coloneqq \frac1{\sqrt{\lambda_1}}
    \left(
        2+
        \frac1{\sqrt{\kappa}}
        +
        \frac1{\kappa}
    \right).
\end{equation}

\item\label{thm:heat-estimates-h1} If \(\Omega\) is convex and
\(
    g\in H_0^1(\Omega),
\)
then
\begin{equation}\label{eq:u-increased-reg}
    u\in L^\infty(0,T;H_0^1(\Omega))
    \cap L^2(0,T;H^2(\Omega)),
    \qquad
    u'\in L^2(0,T;L^2(\Omega)).
\end{equation}
Moreover,
\begin{equation}\label{eq:explicit-H1-H2-estimate}
\begin{split}
    &\operatorname*{ess\,sup}_{0\le t\le T}\norm{u(t)}_{H_0^1(\Omega)}
    +
    \norm{u}_{L^2(0,T;H^2(\Omega))}
    +
    \norm{u'}_{L^2(0,T;L^2(\Omega))}\\
    &\qquad\le
    C_1
    \norm{g}_{H_0^1(\Omega)}
    +
    C_2
    \norm{F}_{L^2(0,T;L^2(\Omega))},
\end{split}
\end{equation}
where \(C_\Omega\) is the constant from \cref{lem:convex-domain-A-H2}, and
\begin{equation}\label{eq:C-H1-H2}
\begin{gathered}
    C_1
    =
    1+\sqrt{\kappa}+\frac{C_\Omega}{\sqrt{\kappa}},
    \quad
    C_2
    =
    2+\frac1{\sqrt{\kappa}}+\frac{C_\Omega}{\kappa}.
\end{gathered}
\end{equation}
\end{enumerate}
\end{theorem}

\begin{proof}[Proof of \cref{thm:heat-estimates}]
The existence and uniqueness statements are standard results, \cite[Section 7.1.2]{evans_2010}. For the estimates we base the proofs on \cite{evans_2010}, but we need to keep track of the constants in order to get explicit bounds, and we avoid the smooth boundary assumption using \cref{lem:convex-domain-A-H2}.

\textbf{Part (1)}.
Test the $m$-th Galerkin equation for \eqref{eq:heat-general} against the corresponding approximation \(u_m(t)\).  Poincare's and Young's
inequalities give
\begin{equation}\label{eq:proof-L2-H1-differential}
    \frac12\frac{d}{dt}\norm{u_m(t)}_{L^2}^2
    +
    \kappa\norm{\nabla u_m(t)}_{L^2}^2
    =
    (F(t),u_m(t))_{L^2}
    \le
    \frac{\kappa}{2}\norm{\nabla u_m(t)}_{L^2}^2
    +
    \frac{1}{2\kappa\lambda_1}\norm{F(t)}_{L^2}^2.
\end{equation}
Integrating over \((0,t)\), taking the supremum over \((0,  T)\), and
passing to the weak limit, lower semicontinuity gives
\begin{equation}\label{eq:proof-L2-energy-bounds}
\begin{split}
    \sup_{0\le t\le T}\norm{u(t)}_{L^2}
    &\le
    \norm{g}_{L^2}
    +
    \frac1{\sqrt{\kappa\lambda_1}}
    \norm{F}_{L^2(0,T;L^2)},
    \\
    \norm{u}_{L^2(0,T;H_0^1)}
    &\le
    \frac1{\sqrt{\kappa}}\norm{g}_{L^2}
    +
    \frac1{\kappa\sqrt{\lambda_1}}
    \norm{F}_{L^2(0,T;L^2)}.
\end{split}
\end{equation}
To bound the third norm, note that
the dual norm satisfies
\begin{equation}\label{eq:proof-Hminus1-dual-bounds}
    \norm{F(t)}_{H^{-1}}
    \le
    \lambda_1^{-1/2}\norm{F(t)}_{L^2},
    \qquad
    \norm{Au(t)}_{H^{-1}}
    =
    \norm{u(t)}_{H_0^1}.
\end{equation}
Thus, since \(u'=F-\kappa Au\), the second inequality in \eqref{eq:proof-L2-energy-bounds} gives
\begin{equation}\label{eq:proof-uprime-Hminus1}
\begin{split}
    \norm{u'}_{L^2(0,T;H^{-1})}
    \le
    \lambda_1^{-1/2}
    \norm{F}_{L^2(0,T;L^2)}
    +
    \kappa
    \norm{u}_{L^2(0,T;H_0^1)}
    \le
    \sqrt{\kappa}\norm{g}_{L^2}
    +
    \frac2{\sqrt{\lambda_1}}
    \norm{F}_{L^2(0,T;L^2)}.
\end{split}
\end{equation}
Combining \eqref{eq:proof-L2-energy-bounds} and \eqref{eq:proof-uprime-Hminus1} proves
\eqref{eq:explicit-L2-estimate}.

\textbf{Part \textup{(2)}}.
Testing the $m$-th Galerkin equation with \(Au_m(t)\) and applying Young's
inequality gives
\begin{equation}\label{eq:proof-H1-Au-identity}
    \frac12\frac{d}{dt}\norm{\nabla u_m(t)}_{L^2}^2
    +
    \kappa\norm{Au_m(t)}_{L^2}^2
    =
    (F(t),Au_m(t))_{L^2}
    \le
    \frac1{2\kappa}\norm{F(t)}_{L^2}^2
    +
    \frac{\kappa}{2}\norm{Au_m(t)}_{L^2}^2.
\end{equation}
Integrating over \((0,t)\), taking the essential supremum over \((0,T)\), and
passing to the weak limit,
\begin{equation}\label{eq:proof-Linf-H01}
    \operatorname*{ess\,sup}_{0\le t\le T}\norm{u(t)}_{H_0^1(\Omega)}
    \le
    \norm{g}_{H_0^1(\Omega)}
    +
    \frac1{\sqrt{\kappa}}
    \norm{F}_{L^2(0,T;L^2(\Omega))},
\end{equation}
and
\begin{equation}\label{eq:proof-Au-L2}
    \norm{Au}_{L^2(0,T;L^2(\Omega))}
    \le
    \frac1{\sqrt{\kappa}}\norm{g}_{H_0^1(\Omega)}
    +
    \frac1{\kappa}
    \norm{F}_{L^2(0,T;L^2(\Omega))}.
\end{equation}
Thus since \(\Omega\) is convex, \cref{lem:convex-domain-A-H2} implies
\begin{equation}\label{eq:proof-H2-via-Au}
    \norm{u}_{L^2(0,T;H^2(\Omega))}
    \le
    C_\Omega\norm{Au}_{L^2(0,T;L^2(\Omega))}
    \le
    C_\Omega
    \left(
        \frac1{\sqrt{\kappa}}\norm{g}_{H_0^1(\Omega)}
        +
        \frac1{\kappa}
        \norm{F}_{L^2(0,T;L^2(\Omega))}
    \right).
\end{equation}
Finally, since \(u'=F-\kappa Au\),
\begin{equation}\label{eq:proof-uprime-L2}
    \norm{u'}_{L^2(0,T;L^2(\Omega))}
    \le
    \sqrt{\kappa}\norm{g}_{H_0^1(\Omega)}
    +
    2\norm{F}_{L^2(0,T;L^2(\Omega))}.
\end{equation}
Adding \eqref{eq:proof-Linf-H01}, \eqref{eq:proof-H2-via-Au}, and
\eqref{eq:proof-uprime-L2} proves \eqref{eq:explicit-H1-H2-estimate}.
\end{proof}

\subsubsection{General wave equation energy estimates}

For the wave equation we provide a single energy estimate matching
\hyperref[thm:heat-estimates-h1]{\cref*{thm:heat-estimates}, part~\textup{(2)}}.

\begin{theorem}\label{thm:wave-estimates}
Let \(\Omega\) be an open, bounded domain and $c, T > 0$. Consider the wave equation
\begin{equation}\label{eq:wave-general}
\begin{cases}
\partial_t^2 u-c^2\Delta u=F & \text{in }\Omega\times(0,T],\\
u=0 & \text{on }\partial\Omega\times[0,T],\\
u=g,\quad \partial_t u=h & \text{on }\Omega\times\set{t=0},
\end{cases}
\end{equation}
where
\[
    F\in L^2(0,T;L^2(\Omega)),
    \qquad
    g\in H_0^1(\Omega),
    \qquad
    h\in L^2(\Omega).
\]
There is a unique weak solution
\[
    u\in L^\infty(0,T;H_0^1(\Omega)),
    \qquad
    u'\in L^\infty(0,T;L^2(\Omega)),
    \qquad
    u''\in L^2(0,T;H^{-1}(\Omega)),
\]
for which the following estimate holds:
\begin{equation}\label{eq:explicit-wave-energy-estimate}
\begin{split}
    &\operatorname*{ess\,sup}_{0<t<T}
    \left(
        \norm{u(t)}_{H_0^1(\Omega)}
        +
        \norm{u'(t)}_{L^2(\Omega)}
    \right)
    +
    \norm{u''}_{L^2(0,T;H^{-1}(\Omega))}
    \\
    &\qquad\le
    \alpha^{\mathrm W}
    \norm{g}_{H_0^1(\Omega)}
    +
    \eta^{\mathrm W}
    \norm{h}_{L^2(\Omega)}
    +
    \beta^{\mathrm W}
    \norm{F}_{L^2(0,T;L^2(\Omega))},
\end{split}
\end{equation}
where
\begin{equation}\label{eq:wave-alpha-beta-constants}
    \alpha^{\mathrm W}\coloneqq 1+c+c^2\sqrt{T},
    \qquad
    \eta^{\mathrm W}\coloneqq 1+\frac1c+c\sqrt{T},
    \qquad
    \beta^{\mathrm W}\coloneqq \frac1{\sqrt{\lambda_1}} + cT
    +\sqrt{T}\left(1+\frac1c\right).
\end{equation}

\end{theorem}

\begin{proof}[Proof of \cref{thm:wave-estimates}]
The existence and uniqueness statements are standard, see
\cite[Section 7.2.2]{evans_2010}. We just need to keep track of constants in the estimates. Testing the $m$-th Galerkin equation with
\(u_m'(t)\) gives
\[
    \frac12\frac{d}{dt}
    \left(
        E_m(t)^2
    \right)
    =
    (F(t),u_m'(t))_{L^2},
\]
where
\[
    E_m(t)^2
    \coloneqq
    \norm{u_m'(t)}_{L^2}^2
    +
    c^2\norm{u_m(t)}_{H_0^1}^2 .
\]
Since \(E_m\) is absolutely continuous, dividing by \(E_m(t)\) where
\(E_m(t)>0\) and applying Cauchy--Schwarz gives, for a.e. \(t\in(0,T)\),
\[
    E_m'(t)
    \le
    \norm{F(t)}_{L^2}.
\]
Thus,
\[
    E_m(t)\le E_m(0)+\int_0^t\norm{F(s)}_{L^2}\,ds .
\]

Taking the essential supremum over \((0,T)\), applying Cauchy--Schwarz to the
resulting integral involving $F$, and passing to the weak limit gives
\begin{equation}\label{eq:proof-wave-energy-bounds}
\begin{split}
    \operatorname*{ess\,sup}_{0<t<T}\norm{u(t)}_{H_0^1}
    &\le
    \norm{g}_{H_0^1}
    +
    \frac1c\norm{h}_{L^2}
    +
    \frac{\sqrt{T}}{c}
    \norm{F}_{L^2(0,T;L^2)},
    \\
    \operatorname*{ess\,sup}_{0<t<T}\norm{u'(t)}_{L^2}
    &\le
    c\norm{g}_{H_0^1}
    +
    \norm{h}_{L^2}
    +
    \sqrt{T}
    \norm{F}_{L^2(0,T;L^2)}.
\end{split}
\end{equation}
To bound the third norm, recall \eqref{eq:proof-Hminus1-dual-bounds}.
Thus, since \(u''=F-c^2Au\), the first inequality in
\eqref{eq:proof-wave-energy-bounds} gives
\begin{equation}\label{eq:proof-wave-utt}
\begin{split}
    \norm{u''}_{L^2(0,T;H^{-1})}
    &\le
    \lambda_1^{-1/2}
    \norm{F}_{L^2(0,T;L^2)}
    +
    c^2
    \norm{u}_{L^2(0,T;H_0^1)}
    \\
    &\le
    c^2\sqrt{T}\norm{g}_{H_0^1}
    +
    c\sqrt{T}\norm{h}_{L^2}
    +
    \left(
        \frac1{\sqrt{\lambda_1}}
        +
        cT
    \right)
    \norm{F}_{L^2(0,T;L^2)}.
\end{split}
\end{equation}
Adding the two inequalities in \eqref{eq:proof-wave-energy-bounds} and
\eqref{eq:proof-wave-utt} proves \eqref{eq:explicit-wave-energy-estimate}.
\end{proof}



\subsubsection{Residual estimates}
Let $u_{\mathrm H}$ be the weak solution of the heat problem in
\eqref{eq:heat-wave-model-framework}, let $u_{\mathrm W}$ be the weak solution
of the wave problem in \eqref{eq:heat-wave-model-framework}, and let
$v \in C^\infty(\Omega)$ vanish on \(\partial\Omega\).
Define
$e_{\mathrm H}\coloneqq u_{\mathrm H}-v$, $e_{\mathrm W}\coloneqq u_{\mathrm W}-v$,
and $v_0, \dot v_0$ as in \eqref{eq:initial-data}. Writing the PDE residuals as
\[
    R_{\mathrm H}(v)\coloneqq \partial_t v-\kappa\Delta v,
    \qquad
    R_{\mathrm W}(v)\coloneqq \partial_t^2v-c^2\Delta v,
\]
then \(e_{\mathrm H}\) and \(e_{\mathrm W}\) are weak solutions of the heat and
wave \emph{residual} equations, respectively,
\begin{equation}\label{eq:heat-wave-residual}
\begin{array}{@{}c@{\qquad}c@{}}
\begin{cases}
\partial_t e-\kappa\Delta e=-R_{\mathrm H}(v) & \text{in }\Omega\times(0,T],\\
e=0 & \text{on }\partial\Omega\times[0,T],\\
e=g-v_0 & \text{on }\Omega\times\set{t=0},
\end{cases}
&
\begin{cases}
\partial_t^2 e-c^2\Delta e=-R_{\mathrm W}(v) & \text{in }\Omega\times(0,T],\\
e=0 & \text{on }\partial\Omega\times[0,T],\\
e=g-v_0,\quad \partial_t e=h-\dot v_0 & \text{on }\Omega\times\set{t=0}.
\end{cases}
\end{array}
\end{equation}

Then the energy estimates from \cref{sec:energy-estimates} give the following corollaries, which are the precise forms of the residual frameworks \eqref{eq:heat-h1-framework}, \eqref{eq:wave-h1-framework} for the heat and wave equations.

\begin{corollary}[Heat residual estimates]\label{cor:heat-residual-estimates}
With $e_H$ defined as above, under the assumptions of
\hyperref[thm:heat-estimates-l2]{\cref*{thm:heat-estimates}, part~\textup{(1)}} and
\hyperref[thm:heat-estimates-h1]{\cref*{thm:heat-estimates}, part~\textup{(2)}},
respectively, it holds that
\begin{subequations}\label{eq:heat-residual-estimates}
\begin{align}
    \sup_{0\le t\le T}\norm{e_{\mathrm H}(t)}_{L^2}
    +\norm{e_{\mathrm H}}_{L^2(0,T;H_0^1)}
    +\norm{\partial_t e_{\mathrm H}}_{L^2(0,T;H^{-1})}
    &\le
    \alpha^{\mathrm H}\norm{g-v_0}_{L^2}
    +
    \beta^{\mathrm H}\norm{R_{\mathrm H}(v)}_{L^2(0,T;L^2)}, \label{eq:heat-residual-estimate-l2}
    \\
    \operatorname*{ess\,sup}_{0<t<T}\norm{e_{\mathrm H}(t)}_{H_0^1}
    +\norm{e_{\mathrm H}}_{L^2(0,T;H^2)}
    +\norm{\partial_t e_{\mathrm H}}_{L^2(0,T;L^2)}
    &\le
    C_1\norm{g-v_0}_{H_0^1}
    +
    C_2\norm{R_{\mathrm H}(v)}_{L^2(0,T;L^2)} . \label{eq:heat-residual-estimate-h1}
\end{align}
\end{subequations}
We denote the left- and right-hand sides of
\eqref{eq:heat-residual-estimate-l2} by
$\mathcal N_{\mathrm H_0}(e_{\mathrm H})$ and
$\mathcal R_{\mathrm H_0}(v)$, respectively, and those of
\eqref{eq:heat-residual-estimate-h1} by
$\mathcal N_{\mathrm H_1}(e_{\mathrm H})$ and
$\mathcal R_{\mathrm H_1}(v)$, respectively.
\end{corollary}

\begin{corollary}[Wave residual estimate]\label{cor:wave-residual-estimate}
With $e_W$ defined as above, under the assumptions of \cref{thm:wave-estimates},
\begin{equation}\label{eq:wave-residual-estimate}
\begin{split}
    &\operatorname*{ess\,sup}_{0<t<T}
    \left(
        \norm{e_{\mathrm W}(t)}_{H_0^1}
        +
        \norm{\partial_t e_{\mathrm W}(t)}_{L^2}
    \right)
    +
    \norm{\partial_t^2 e_{\mathrm W}}_{L^2(0,T;H^{-1})}
    \\
    &\qquad\le
    \alpha^{\mathrm W}\norm{g-v_0}_{H_0^1}
    +
    \eta^{\mathrm W}\norm{h-\dot v_0}_{L^2}
    +
    \beta^{\mathrm W}\norm{R_{\mathrm W}(v)}_{L^2(0,T;L^2)} .
\end{split}
\end{equation}
We denote the left- and right-hand sides of
\eqref{eq:wave-residual-estimate} by
$\mathcal N_{\mathrm W}(e_{\mathrm W})$ and
$\mathcal R_{\mathrm W}(v)$, respectively.
\end{corollary}

Thus the quadrature methods outlined in \cref{subsec:quadrature} must be applied to the residual norms:
\begin{equation}\label{eq:quadrature-target-residual-norms}
\begin{gathered}
    \norm{g-v_0}_{L^2(\Omega)},
    \quad
    \norm{\nabla(g-v_0)}_{L^2(\Omega)},
    \quad
    \norm{h-\dot v_0}_{L^2(\Omega)},
    \\
    \norm{R_{\mathrm H}(v)}_{L^2(0,T;L^2(\Omega))},
    \quad
    \norm{R_{\mathrm W}(v)}_{L^2(0,T;L^2(\Omega))}.
\end{gathered}
\end{equation}


\subsection{Rigorous quadrature for residual norms}\label{sec:residual-quad}
We now describe in detail how the general quadrature framework in \cref{subsec:quadrature} applies to the residual norms.
For concreteness, assume that \(\Omega=(0,1)^d\), $d\in\N$, and that the spatial radius vector $\boldsymbol{\eps}_x\in(0,\infty)^d$ and time radius $\eps_t>0$ satisfy \cref{assumption:rectangular-domain}.  Set $\boldsymbol{\eps}\coloneqq(\boldsymbol{\eps}_x,\eps_t)$.  In particular, we have the partitions
 \(
    \mathcal P_\Omega
    \coloneqq
    \set*{\mathcal B_{\boldsymbol{\eps}_x}(y_k)}_k,
  \) and
  \(  \mathcal P
    \coloneqq
    \set*{\mathcal B_{\boldsymbol{\eps}}(y_k,t_m)}_{k,m}
  \)
  of $\Omega$ and $\Omega_T$, respectively.
In this case the boundary function has the explicit form
\[
B(x) = \prod_{i=1}^d x_i(1-x_i),
\]
whose explicit moduli $\omega_B^\beta$ satisfying \cref{ass:B-moduli} are recorded in
\cref{app:boundary-factor-moduli}. Since the residual estimates and their bounds are independent of training methodology, we need only assume that the approximation is given by $v = Bf$ with $f \in \cF_{L, w}$ for some fixed $L, w \in \N$.




Let $\phi$ be one of the five integrands in
\eqref{eq:quadrature-target-residual-norms}.  For every cell, the construction
in \cref{subsubsec:cellwise-quadrature-bounds} requires only verified first-
or second-derivative coefficients as in
\eqref{eq:cell-derivative-coefficients}; its pointwise envelopes and all
cross and squared-remainder terms are then integrated analytically through
the moments \eqref{eq:centered-cell-moments}.  The neural-network derivative
bounds entering these coefficients are computed by
\cref{alg:first-deriv-bounds,alg:higher-deriv-bounds}.

For the initial-data residuals, let
$C_k=\mathcal B_{\boldsymbol{\eps}_x}(y_k)\in\mathcal P_\Omega$ and set
$\boldsymbol{\eps}_0\coloneqq(\boldsymbol{\eps}_x,0)$.  Network bounds are
centered at $(y_k,0)$ with radius vector $\boldsymbol{\eps}_0$, and hence hold
on the initial slice
\[
 \mathcal B_{\boldsymbol{\eps}_0}(y_k,0)
 =C_k\times\{0\}.
\]
The complete product-rule formulas for all initial-data coefficients are
given in \cref{app:initial-data-derivative-coefficients}.

\subsubsection{Initial displacement residuals}\label{sec:initial-displacement-quad}

Set $e_0\coloneqq g-v_0=g-Bf(\cdot,0)$.  For $Q^0$ we use the coefficients
$b_{q,C_k}(e_0)$ in
\eqref{eq:appendix-initial-displacement-first-coefficients}; this requires
point values and a Lipschitz constant for $g$.  For $Q^1$ we use
$h_{q\ell,C_k}(e_0)$ in
\eqref{eq:appendix-initial-displacement-second-coefficients}; the point values
and Lipschitz constants of the first derivatives of $g$ also make the affine
model computable.  In both cases, inserting these coefficients into
\cref{eq:cell-remainder-envelopes,eq:q0-cell-moment-correction,eq:q1-cell-moment-correction}
and summing over the spatial partition gives
\begin{equation}\label{eq:initial-displacement-global-bound}
    \norm{e_0}_{L^2(\Omega)}
    \leq
    \left(
        Q_{\mathcal P_\Omega}^r(e_0)
        +E_{\mathcal P_\Omega}^r(e_0)
    \right)^{1/2}
    \eqqcolon \mathcal{U}^r_{\mathcal{P}_\Omega}(e_0),
    \qquad r\in\{0,1\}.
\end{equation}

\subsubsection{Initial-displacement gradient residual}\label{sec:initial-displacement-gradient-quad}

By \cref{rem:vector-residual-bounds}, it suffices to bound each
component $\partial_j e_0$, $j=1,\dots,d$.  Its $Q^0$ and $Q^1$
coefficients are, respectively,
\eqref{eq:appendix-initial-gradient-first-coefficients} and
\eqref{eq:appendix-initial-gradient-second-coefficients}.  The latter requires
point values and Lipschitz constants through the second spatial derivatives
of $g$.  Applying the common moment construction componentwise gives
\begin{equation}\label{eq:initial-gradient-global-bound}
    \norm{\nabla e_0}_{L^2(\Omega)}
    \leq
    \left(
        \sum_{j=1}^d
        \left(Q_{\mathcal P_\Omega}^r(\partial_j e_0)
        +E_{\mathcal P_\Omega}^r(\partial_j e_0)\right)
    \right)^{1/2}
    \eqqcolon \mathcal U_{\mathcal P_\Omega}^r(\nabla e_0),
    \qquad r\in\{0,1\}.
\end{equation}

\subsubsection{Initial velocity residual} \label{sec:initial-velocity-quad}

For the initial velocity residual, consider
\[
    e_{\mathrm v}\coloneqq h-\dot v_0
    =h-B\partial_t f(\cdot,0).
\]
The same construction applies with $g$ replaced by $h$ and $f(\cdot,0)$ by
$\partial_tf(\cdot,0)$.  Specifically, $Q^0$ uses
\eqref{eq:appendix-initial-velocity-first-coefficients}, while $Q^1$ uses
\eqref{eq:appendix-initial-velocity-second-coefficients} together with the
center value and spatial gradient of $e_{\mathrm v}$.  Hence
\begin{equation}\label{eq:initial-velocity-global-bound}
    \norm{e_{\mathrm v}}_{L^2(\Omega)}
    \leq
    \left(
        Q_{\mathcal P_\Omega}^r(e_{\mathrm v})
        +E_{\mathcal P_\Omega}^r(e_{\mathrm v})
    \right)^{1/2}
    \eqqcolon \mathcal U_{\mathcal P_\Omega}^r(e_{\mathrm v}),
    \qquad r\in\{0,1\}.
\end{equation}

\subsubsection{PDE residuals}\label{sec:pde-residual-quad}

The PDE residuals live on spacetime and involve second derivatives of the
network, so they are the most demanding terms numerically.  We use the affine
rule on a spacetime partition $\mathcal T$ with anisotropic cells
$C=\zeta_C+\prod_{q=1}^{d+1}[-\eps_{C,q},\eps_{C,q}]$, allowing adaptive
refinement in high dimensions.  For
$\mathfrak p\in\{\mathrm H,\mathrm W\}$, the coefficients
\begin{equation}\label{eq:pde-residual-second-coefficients}
 h_{q\ell,C}(R_{\mathfrak p}(Bf))
 \coloneqq h_{q\ell,C}^{\mathfrak p}
 \geq\sup_{z\in C}|\partial_{q\ell}R_{\mathfrak p}(Bf)(z)|,
 \qquad q,\ell=1,\dots,d+1,
\end{equation}
are computed in \cref{app:pde-residual-hessian-bounds} using network
derivatives through order four.  Inserting them into the same affine
remainder envelope and moment correction
\cref{eq:cell-remainder-envelopes,eq:q1-cell-moment-correction} gives
\begin{equation}\label{eq:pde-global-bound}
 \norm{R_{\mathfrak p}(Bf)}_{L^2(\Omega_T)}
 \leq
 \left(Q_{\mathcal T}^1(R_{\mathfrak p}(Bf))
 +E_{\mathcal T}^1(R_{\mathfrak p}(Bf))\right)^{1/2}
 \eqqcolon\mathcal U_{\mathcal T}^1(R_{\mathfrak p}(v)).
\end{equation}

\subsection{Verification algorithms}
\label{subsec:verification-algorithms}

We give an informal \emph{local} verification algorithm and a \emph{global} one
which together affirmatively answer \cref{question-sequence}, illustrating the residual verification framework for our model problems.
For spatial and spacetime partitions $\mathcal T_\Omega$ and $\mathcal T$,
respectively, the preceding constructions give computable bounds for the quadrature rules $r\in\{0,1\}$, which we denote by
\begin{equation}\label{eq:computable-residual-bounds}
\begin{aligned}
    \widehat{\mathcal R}_{\mathrm H_0}^{r}
    (v;\mathcal T_\Omega,\mathcal T)
    &\coloneqq
    \alpha^{\mathrm H}\mathcal U_{\mathcal T_\Omega}^r(e_0)
    +\beta^{\mathrm H}\mathcal U_{\mathcal T}^1(R_{\mathrm H}(v)),
    \\
    \widehat{\mathcal R}_{\mathrm H_1}^{r}
    (v;\mathcal T_\Omega,\mathcal T)
    &\coloneqq
    C_1\mathcal U_{\mathcal T_\Omega}^r(\nabla e_0)
    +C_2\mathcal U_{\mathcal T}^1(R_{\mathrm H}(v)),    \\
    \widehat{\mathcal R}_{\mathrm W}^r
    (v;\mathcal T_\Omega,\mathcal T)
    &\coloneqq
    \alpha^{\mathrm W}\mathcal U_{\mathcal T_\Omega}^r(\nabla e_0)
    +\eta^{\mathrm W}\mathcal U_{\mathcal T_\Omega}^r(e_{\mathrm v})
    +\beta^{\mathrm W}\mathcal U_{\mathcal T}^1
        (R_{\mathrm W}(v)). 
\end{aligned}
\end{equation}
Then
\cref{cor:heat-residual-estimates,cor:wave-residual-estimate} and the
quadrature constructions above give
\begin{equation}\label{eq:residual-bound-chain}
    \mathcal N_\sigma(u_\sigma-v)
    \leq \mathcal R_\sigma(v)
    \leq \widehat{\mathcal R}_\sigma^r
    (v;\mathcal T_\Omega,\mathcal T), \qquad \sigma\in\{\mathrm H_0,\mathrm H_1,\mathrm W\},
\end{equation}
where $u_{\mathrm H_0}=u_{\mathrm H_1}=u_{\mathrm H}$.

\begin{theorem}[Convergence of computable residual bounds]
\label{thm:computable-residual-bound-convergence}
Fix
$v=Bf$ as in \eqref{eq:boundary-ansatz} and suppose that $\Omega_T$ is covered by partitions $\mathcal T_\Omega$ and $\mathcal T$ as in \cref{assumption:rectangular-domain}.  Let
$\sigma\in\{\mathrm H_0,\mathrm H_1,\mathrm W\}$, and $r\in\{0,1\}$.
Let $h_\Omega$ and $h_T$ be the maximal coordinate radii
of the cells in $\mathcal T_\Omega$ and $\mathcal T$, respectively.  There
is a constant $K>0$, independent of the partitions, such that
\begin{equation}\label{eq:computable-residual-bound-convergence}
    0\leq
    \widehat{\mathcal R}_{\sigma}^r
    (v;\mathcal T_\Omega,\mathcal T)-\mathcal R_{\sigma}(v)
    \leq K\left(h_\Omega^{r+1}+h_T^2\right).
\end{equation}
Consequently, as $h_\Omega,h_T\to0$,
\[
    \widehat{\mathcal R}_{\sigma}^r
    (v;\mathcal T_\Omega,\mathcal T)
    \longrightarrow \mathcal R_{\sigma}(v).
\]
\end{theorem}

\begin{proof}
For any scalar residual integrand $\phi$, let $P_{\mathcal T}^r$ and
$D_{\mathcal T}^r(\phi)$ denote the piecewise polynomial and remainder
envelope from \cref{eq:cell-polynomial-models,eq:cell-remainder-envelopes},
and set
\[
    \mathfrak D_{\mathcal T}^r(\phi)
    \coloneqq
    \left(\sum_{C\in\mathcal T}
    \norm{D_C^r(\phi)}_{L^2(C)}^2\right)^{1/2}.
\]
For $s$-dimensional cells, the definitions of $A_C^r$ and $Q_C^r$ imply
\[
 \norm{A_C^r(\phi)}_{L^2(C)}
 \leq c_s Q_C^r(\phi)^{1/2},
 \qquad c_s\coloneqq\sqrt{s+1},
\]
where for $r=0$ one may take $c_s=1$.  Cauchy--Schwarz over the cells,
\eqref{eq:cell-verified-quadrature-error}, and the reverse triangle
inequality therefore give
\[
 \begin{aligned}
  \mathcal U_{\mathcal T}^r(\phi)
  &\leq Q_{\mathcal T}^r(\phi)^{1/2}
  +c_s\mathfrak D_{\mathcal T}^r(\phi),\\
  0\leq\mathcal U_{\mathcal T}^r(\phi)-\norm{\phi}_{L^2}
  &\leq(c_s+1)\mathfrak D_{\mathcal T}^r(\phi).
 \end{aligned}
\]
The uniform derivative coefficients in
\eqref{eq:cell-remainder-envelopes} show that
$\mathfrak D_{\mathcal T_\Omega}^r(\phi)
\leq C_\phi|\Omega|^{1/2}h_\Omega^{r+1}$ for every scalar initial-data
integrand.  The same componentwise estimate, followed by the Euclidean
triangle inequality, applies to $\nabla e_0$.  For either PDE residual, its
bounded Hessian gives
$\mathfrak D_{\mathcal T}^1(R_{\mathfrak p}(v))
\leq C_{\mathfrak p}|\Omega_T|^{1/2}h_T^2$.
Substitution into \eqref{eq:computable-residual-bounds}, together with the
second inequality in \eqref{eq:residual-bound-chain}, proves
\eqref{eq:computable-residual-bound-convergence}.
\end{proof}


\subsubsection{Local verification algorithm}

Given $v=Bf$, an estimate
$\sigma\in\{\mathrm H_0,\mathrm H_1,\mathrm W\}$, and a desired accuracy
$\eps_0>0$, we want to verify that $v$ is within $\eps_0$ of the unknown
solution $u_\sigma$ in the corresponding error norm.

\begin{enumerate}
    \item input: $v=Bf$, $\sigma\in\{\mathrm H_0,\mathrm H_1,\mathrm W\}$,
    $\eps_0>0$, the data point values and Lipschitz constants, and the
    boundary-factor moduli from
    \cref{ass:data-moduli,ass:B-moduli};
    \item initialize a spatial radius vector $\boldsymbol{\eps}_x\in(0,\infty)^d$
    and a time radius $\eps_t>0$ for which the spatial and spacetime
    partitions $\mathcal T_\Omega$ and $\mathcal T$ exist as in
    \cref{assumption:rectangular-domain};
    \item compute the initial-data and PDE-residual norm bounds over
    these partitions using the algorithms from \cref{sec:activation} and the
    constructions above, and combine into
    $\widehat{\mathcal R}_{\sigma}^r
    (v;\mathcal T_\Omega,\mathcal T)$;
    \item if
    \begin{equation}\label{eq:ima-halting-criterion}
        \widehat{\mathcal R}_{\sigma}^r
        (v;\mathcal T_\Omega,\mathcal T)<\eps_0,
    \end{equation}
    halt and return \textsc{yes}. Otherwise replace
    $(\boldsymbol{\eps}_x,\eps_t)$ by
    $(\boldsymbol{\eps}_x/2,\eps_t/2)$, refine the partitions accordingly,
    and start over from point~3.
\end{enumerate}

Even when the error of $v$ is smaller than $\eps_0$, there is in general no
guarantee that the local verification algorithm will ever halt, since the
residual estimate itself may be pessimistic. This motivates the global
algorithm.

\subsubsection{Global verification algorithm}

Fix $\sigma\in\{\mathrm H_0,\mathrm H_1,\mathrm W\}$.  Repeatedly using the local algorithm, we can formulate the
global algorithm which affirmatively answers \cref{question-sequence}. The explicit
construction above applies to the heat and wave equations, and the same
procedure works for any PDE admitting a residual framework of the form
\eqref{eq:XYZ-residuals} whose residual norms can be bounded by our
quadrature and derivative-bound algorithms.

\begin{enumerate}
    \item input: a tolerance $\eps_0>0$, the estimate $\sigma$, and a sequence
    $\{v_m=Bf_m\}$ containing a subsequence $\{v_{m_k}\}$ converging to $u$
    in a space $Z$ for which
    \eqref{eq:XYZ-residuals} holds.
    \item iteratively start running the local algorithm on each triple
    $(v_m,\sigma,\eps_0)$.
    \item output the first index $m$ for which the local algorithm halts.
\end{enumerate}

The global algorithm halts after finitely many steps. Indeed, by \eqref{eq:XYZ-residuals} there is an index $m_k$ for which
$\mathcal R_\sigma(v_{m_k})<\eps_0$.  For this fixed $v_{m_k}$,
\cref{thm:computable-residual-bound-convergence} implies that after finitely
many refinements the local algorithm satisfies
\eqref{eq:ima-halting-criterion}. Thus the global algorithm halts and returns
an approximation $v_m$ satisfying $\norm{u-v_m}_X<\eps_0$.
